\documentclass[10pt,a4paper]{amsart}
\usepackage[margin=19mm]{geometry}
\usepackage{amsmath,amssymb,mathtools}
\usepackage[scr=rsfs,cal=euler]{mathalfa}
\usepackage{xfrac}
\usepackage[unicode,hidelinks,bookmarks=false]{hyperref}
\newcommand{\ie}{i.e.\ }
\newcommand{\cf}{cf.\ }
\newcommand{\resp}{respectively\ }
\newcommand{\Subsection}{\subsection}
\providecommand{\nobreakdash}{\nobreak\hbox{-}}
\setlength{\emergencystretch}{2em}
\multlinegap0pt
\usepackage{bm}
\usepackage{bbm}
\usepackage{dsfont}
\usepackage{xspace}
\usepackage{cancel}
\usepackage{mathtools}
\usepackage{amsthm}
\makeatletter
\@ifundefined{theorem}{\newtheorem{theorem}{Theorem}}{}
\makeatother
\title[Rational function approximation with normalized positive den]{Rational function approximation with normalized positive denominators}
\author{James Chok, Geoffrey M. Vasil}
\date{Source: arXiv:2310.12053v4 (CC BY 4.0)}
\begin{document}
\maketitle

\noindent\emph{Source and licence.} Rational function approximation with normalized positive denominators. Version arXiv:2310.12053v4 (CC BY 4.0), distributed under the Creative Commons Attribution 4.0 International licence, as recorded in the arXiv licence metadata for this exact version and shown on the abstract page https://arxiv.org/abs/2310.12053v4. Full rights notice: licenses/arxiv-2310.12053-CC-BY-4.0.txt.\par\smallskip

\noindent\textit{Editorial scope.} Counted unit: the Theorem whose source label is \texttt{None} in the article (source0.tex lines 244--249, source label \texttt{None}), delivered together with the article's own complete original proof of it (source0.tex lines 250--272, one \texttt{proof} environment of 23 lines). No mathematics is written, repaired or re-derived: statement and proof are the article's own text line for line. Required context retained uncounted: none. Every definition, display and label of those lines is retained unchanged. Omitted from this package and not redistributed: 1--243, 273--976. Nothing inside a retained excerpt is omitted or altered: every retained body is delivered exactly as the article's lines are written, so no omission receipt of any kind accompanies this package. Citations: the retained excerpts carry no citation command at all, so no bibliography entry of the article is transcribed and no reference is redistributed. Reference adaptation: none was needed and none was made; every reference inside the delivered unit resolves inside it, so no unresolved reference or citation is produced. Printed slips kept literal: no printed word, symbol, hypothesis or number of the counted statement or of its proof was altered. Layout and class: the article's own class is \texttt{article} with packages including arxiv, inputenc, fontenc, booktabs, nicefrac, microtype, lipsum, doi, hyperref, url, amsfonts, amsmath, amssymb, parskip; the wrapper is the lane's pinned \texttt{amsart} editorial wrapper at 10pt on A4, which restates the theorem-like environments the retained text needs and adds the article's own macro definitions that the retained text needs (none), with extra wrapper packages (bm, bbm, dsfont, xspace, cancel, mathtools). The delivered numbering is the unit's own. Assets: no retained span contains a figure environment, a graphics-inclusion command, a PDF-page inclusion, a table or a TikZ picture; no image file is copied into this package, the package declares no asset dependency and no image file is redistributed with it. Licence: version arXiv:2310.12053v4 of this article is distributed under the Creative Commons Attribution 4.0 International licence, as recorded in the arXiv licence metadata for this exact version and shown on the abstract page https://arxiv.org/abs/2310.12053v4. A full rights notice is delivered at licenses/arxiv-2310.12053-CC-BY-4.0.txt. Characters: every retained excerpt is pure ASCII, so the delivered engine, XeLaTeX with its default Latin Modern text font, reports no missing character.
\begin{theorem}
Any polynomial in $\overline{\mathcal{P}}(n)$ has an equivalent Bernstein basis representation with coefficients that are all positive, i.e., there exist Bernstein coefficients $b_i\geq 0$ such that
\begin{align}
    p(x)\ :=\ 1 + a_1x + \cdots + a_n x^n\ =\ \sum_{i=0}^n\, b_i\, \mathcal{B}^{(n)}_i(x), \quad\text{with}\quad \sum_{i=1}^n\, |a_i|\ =\ 1.
\end{align}
\end{theorem}

\begin{proof}
Since $p\in \overline{\mathcal{P}}(n)$, then $\sum_i|a_i|=1$, and hence
\begin{align}
    p(x)\ &=\ \sum_{a_i\geq 0}\,|a_i|\,(1 + x^i)\ +\ \sum_{a_i < 0}\,|a_i|\,(1 - x^i).
\end{align}
Thus, to account for the sign, these polynomials can be written as a convex combination of
\begin{align}
    M_{i, \pm}(x)\ =\ 1 \pm x^i,\quad \text{for}\quad 1\leq i\leq n,
\end{align}
with positive weight given to one of $M_{i,\pm}$ with the correct sign and zero to the other. That is, there exists positive weights, $w_{i, s}\geq 0$, such that 
\begin{align}
    p(x)\ =\ \sum_{s\in\{+, -\}}\,\sum_{i=1}^n\, w_{i, s}\, M_{i, s}(x)\quad \text{and}\quad \sum_{s\in\{+, -\}}\, \sum_{i=1}^n\, w_{i, s}\ =\ 1.
\end{align}

From the Binomial and Vandermonde's identity, the monomial coefficients, $a_i$, can be expressed in Bernstein coefficients via $b_j = \sum_{k=0}^j a_k \binom{j}{k} / \binom{n}{k}$. Thus
\begin{align}
    M_{i, \pm}(x)\ =\ 1 \pm x^i\ =\ \sum_{j}\, b_j\, \mathcal{B}^{(n)}_j(x)\quad\text{where}\quad b_j\ =\ 1 \pm \frac{\binom{j}{i}}{\binom{n}{i}}\mathbbm{1}_{i\leq j},
\end{align}
and $\mathbbm{1}_{i\leq j} = 1$ if $i\leq j$ and $0$ otherwise. It is clear that $b_j\geq 0$ for $M_{i, +}$. For $M_{i, -}$, since $j\leq n$, it follows from the geometric interpretation of combinations that $\binom{j}{i}\leq\binom{n}{i}$.
Thus $b_j\geq 0$ for $M_{i, \pm}$. Hence, each basis function, $M_{i, \pm}$, represents a positively weighted sum of Bernstein polynomials.

Since every polynomial in $\overline{\mathcal{P}}(n)$ can be written as a convex combination of $M_{i, \pm}$, it follows that it also can be written as a convex combination of Bernstein polynomials with positive coefficients. Thus, every polynomial in $\overline{\mathcal{P}}(n)$ can be written as a sum of Bernstein polynomials with positive coefficients.
\end{proof}

\end{document}
